\startfirstchapter{Introduction}
\label{chapter:introduction}
Canada has a growing problem of abandoned oil wells and contaminated petrochemical legacy silos. Surveys estimate that 15,682 abandoned wells require plugging and/or land reclamation in Alberta and Saskatchewan alone \cite{est_cost_of_cleaning_canada}.  The presence of petroleum hydrocarbons (PHC) in soil poses serious risks to the environment, agricultural productivity, and human health in surrounding communities \cite{PHCBadForCrops,PHCBadForHumans,PHCBadForEnv}. As the world transitions away from fossil fuels, more facilities will be decommissioned under inadequate closure procedures, and the scale of this problem will only increase.

In situ remediation technologies artificially accelerate biodegradation within the soil without excavating the contaminated material \cite{remTechReview}. Biostimulation, one such proven approach, introduces nutrients directly into PHC-contaminated sites to promote microbial activity and hasten the remediation process. Compared to ex situ methods which require physically removing and treating soil off-site; in situ techniques are substantially less costly and less environmentally disruptive \cite{LioraTech}. However, they demand a much richer understanding of soil conditions to treat contaminated zones accurately, and they rely on measurements to verify treatment progress. Historically, these measurements were collected by hand, yielding data that was sparse both spatially and temporally. In recent years, IoT sensor networks deployed at contaminated sites have enabled continuous monitoring of pollutant concentrations and environmental conditions \cite{LioraRemoteSensing,atmMeasInSoil}, providing real-time estimates of how a PHC plume is distributed and evolving over time.

This adoption of continuous remote sensing has largely resolved the temporal sparsity problem. Spatial sparsity, however, remains a significant challenge. Deploying and installing dense sensor arrays is expensive, leaving large gaps in coverage. Without an accurate spatial map of the PHC plume, remediators cannot effectively target treatment, resulting in wasted resources and land that remains unusable long after intervention.

Spatial interpolation is the logical method to reconstruct plume distribution from sparse measurements. However, these methods originate from a variety of fields such as geostatistics, numerical weather prediction and inverse modelling which makes systematic comparison difficult. This thesis addresses that gap by evaluating several interpolation methods on both synthetic and real-world data to identify the most effective method for petrochemical remediation in Canada. 

The spatial interpolation methods that will be compared are Inverse Distance Weighting (IDW), Kriging, 4D-Variational Data Assimilation (4D-Var), the Gauss-Levenberg-Marquardt algorithm (GLM) and the Iterative Ensemble Smoother (IES). These methods can be grouped into two categories. IDW and Kriging are purely distance-based interpolators that estimate the concentration at unsampled locations as a weighted sum of surrounding measurements using a chosen model. By contrast, 4D-Var, GLM and IES are parameter estimation algorithms that fit a given physics-based contaminant transport model to temporal measurements. Once calibrated, the model provides PHC concentration estimates at any unmeasured locations within and around the plume. 

Each method is evaluated on two datasets. A synthetic soil contaminant dataset with a dense ground truth is used to tune hyperparameters. Additionally, robustness to measurement noise, soil heterogeneity, sensor dropout, and spatial and temporal sparsity will be assessed with modifications to the synthetic data. Building on this controlled evaluation, the methods are then applied to real sensor data collected in the Canadian prairies using leave-$n$-out cross-validation.



\textcolor{red}{Need a summary of the results and a roadmap for the sections here still} 


% \input chapters/1/intro_how_to